\section{Identifiability-Aware Apportionment Method}
\label{sec:identifiability_aware_method}

We combine the estimator from Section~\ref{sec:wind_conditioned_source_apportionment}
with the diagnostics from Section~\ref{sec:identifiability_theory}.  The result
is an identifiability-aware source apportionment procedure, abbreviated IASA,
that estimates source contributions only at the resolution supported by the
projected response matrix.

\subsection{Inputs and Outputs}
\label{subsec:iasa_inputs_outputs}

IASA takes sparse pollution observations and their row mask, sensor locations
\(X_{\mathrm{sens}}\), wind observations and masks, source
inventories \(S\in\mathbb R^{q\times K}\), a background basis \(Q\), a maximum
lag \(L\), a nonnegative temporal basis \(\Phi\in\mathbb R_+^{T\times B}\),
and an open-boundary transport response \(G^\partial\).  It also uses a
numerical SVD tolerance \(\tau_{\mathrm{num}}\) and distinct thresholds
\(\tau_\sigma\), \(\tau_\rho\), and \(\tau_v\) for effective rank, pairwise
ambiguity, and visibility.

The outputs are source--basis coefficient estimates, reconstructed source
activities and contribution estimates at the identifiable grouping, fitted
trajectories and residuals, uncertainty intervals, singular-value, coherence,
and ray-distance diagnostics, low-visibility flags, and merge recommendations.

The thresholds are summarized in Table~\ref{tab:thresholds}; each is predeclared
or calibrated from noise or physics and none is tuned on source-recovery results.

\begin{table}[t]
\centering
\small
\setlength{\tabcolsep}{3pt}
\caption{IASA thresholds, their meaning, and how they are set.  None is
optimized on source-recovery error.}
\label{tab:thresholds}
\begin{tabular}{lll}
\toprule
Threshold & Meaning & How set \\
\midrule
\(\tau_{\mathrm{num}}\) & numerical rank tol.\ & \(\max(N,J)\epsilon_{\mathrm{mach}}\sigma_1\) \\
\(\tau_\sigma\) & effective-rank floor & noise scale, e.g.\ \(\sigma_e\sqrt N\) \\
\(\tau_v\) & visibility floor & \(\mathrm{SNR}_{\min}\,\sigma_e\) \\
\(\tau_\rho\) & ambiguity coherence & predeclared, near \(1\) \\
\(\tau_\rho^{\mathrm{ref}}\) & refinement coherence & default \(\tau_\rho^{\mathrm{ref}}=\tau_\rho\) \\
\(\epsilon_w\) & wind-drift cap & physical wind tolerance \\
\bottomrule
\end{tabular}
\end{table}

\subsection{Procedure}
\label{subsec:iasa_procedure}

Algorithm~\ref{alg:iasa} states the procedure.  Its step order enforces the
method's central discipline of separating fit from diagnosis: all model choices,
thresholds, and the fixed-zero mask are predeclared (line~1), the response is
built, projected, and fitted (lines~2--5), and only afterwards are the
identifiability diagnostics computed and conservative groups reported
(lines~6--8).  Because no threshold or column set is chosen after inspecting the
fitted coefficients---the lag itself is selected from response-matrix geometry,
not from \(\widehat{\mathbf c}\)---the reported identifiability cannot be made
artificially favorable by post-fit tuning.

\begin{algorithm}[t]
\caption{Identifiability-Aware Source Apportionment (IASA).}
\label{alg:iasa}
\footnotesize
\textbf{Input:} observations \(Y\) and observed-row mask, sensors
\(X_{\mathrm{sens}}\), wind observations and masks, inventories
\(S\in\mathbb R^{q\times K}\), background basis \(Q\), maximum lag \(L\),
temporal basis \(\Phi\in\mathbb R_+^{T\times B}\), transport response
\(G^\partial\), thresholds \(\tau_{\mathrm{num}},\tau_\sigma,\tau_\rho,\tau_v\).\\
\textbf{Output:} source--basis coefficients and activities at the identifiable
grouping, fitted trajectories/residuals, uncertainty intervals, diagnostics,
low-visibility flags, and merge recommendations.
\begin{algsteps}
  \item \textbf{Predeclare} the primary \(Q\), the physical lag-candidate grid,
  \(\Phi\), the inventory version, and the fixed-zero mask \(\mathcal F_0\),
  before inspecting any fit.
  \item \textbf{Wind:} convert \(WD/WS\) to transport vectors
  (Eq.~\ref{eq:wind_direction_conversion}), preserving masks/observed values for
  audit; query the adopted kernel imputer on each response-grid cell for
  \(\widehat{\mathbf w}_t(\mathbf x_g)\) and convert to grid displacement.
  \item \textbf{Response \& lag:} over unmasked source--basis pairs build
  \(\mathbf h_{t,kb}^{\mathrm{lag}}\) (Eq.~\ref{eq:constructed_response}); select
  the smallest \(L\) with \(\eta_L\le\tau_L\) from geometry alone
  (Eq.~\ref{eq:lag_convergence}), retaining the full lag sweep.
  \item \textbf{Align \& project:} apply the observed-PM\(_{2.5}\) selection
  identically to \(Y,H,Q\) and metadata (Eq.~\ref{eq:observation_selection}) with
  source-major, basis-minor columns; form \(\widetilde Y,\widetilde H_\Phi\)
  (Eq.~\ref{eq:projected_model}).
  \item \textbf{Fit:} solve \(\widehat{\mathbf c}\)
  (Eq.~\ref{eq:nnls_estimator}) and reconstruct \(\widehat\theta_k\)
  (Eq.~\ref{eq:reconstructed_activity}).
  \item \textbf{Diagnose:} compute \(r_{\mathrm{num}}\), the singular values,
  \(r_{\mathrm{eff}}(\tau_\sigma)\), visibility \(v_j\), absorption \(a_j\),
  coherence \(\rho_{ij}\), and ray distance \(d^{\mathrm{ray}}_{ij}\).
  \item \textbf{Flag:} weak set \(\mathcal W=\{j:v_j\le\tau_v\}\); eligible
  ambiguous pairs \(\mathcal A=\{(i,j):i,j\notin\mathcal W,\rho_{ij}>\tau_\rho\}\).
  \item \textbf{Group:} take deterministic connected components of the source
  ambiguity graph as conservative report groups
  (Sec.~\ref{subsec:source_merging}); weak coefficients are flagged but create no
  edges, and group activities/contributions are member sums that do not replace
  the coefficient fit.
  \item \textbf{Adequacy:} if an externally calibrated noise model exists, run
  the refitted parametric-bootstrap residual test
  (Sec.~\ref{subsec:model_adequacy}); otherwise label the residual summary
  uncalibrated.
  \item \textbf{Report} source estimates, observation/transport uncertainty,
  inventory robustness scenarios, residual adequacy, diagnostics, and
  conservative source groups.
\end{algsteps}
\end{algorithm}

This procedure separates two questions: which source activity trajectories best
fit the projected sensor data, and which source--basis coefficients are
supported by the geometry of \(\widetilde H_\Phi\).

The objective in Equation~\ref{eq:nnls_estimator} is a smooth convex quadratic
under a simple nonnegativity constraint, so we solve it with projected FISTA
\citep{beck2009fast}: an accelerated proximal-gradient method whose proximal step
for the indicator of \(\{\mathbf c\ge0\}\) is the clamp
\(\mathbf c\leftarrow\max(\mathbf c,0)\).  Each gradient step uses a Lipschitz
step size derived from the largest singular value, with monotone restart, and
convergence is checked on the projected-gradient/KKT residual and the relative
objective change.  Device, dtype, convergence status, iteration count,
    and the final KKT residual are part of the fit record.  Transport-ensemble
    and bootstrap refits are batched on the same device when their shapes agree.

\subsection{Uncertainty Reporting}
\label{subsec:uncertainty_reporting}
Let
\begin{equation}
\widetilde r=\widetilde Y-\widetilde H_\Phi\widehat{\mathbf c}
\label{eq:projected_residual}
\end{equation}
be the projected residual.  A residual variance estimate is
\begin{equation}
\widehat\sigma^2
=
\frac{\|\widetilde r\|_2^2}{\max(N-r_{\mathrm{eff}}(\tau_\sigma),1)}.
\label{eq:residual_variance}
\end{equation}
For an unconstrained or active-set approximation, let
\(\mathcal A_c=\{j:\widehat c_j>\tau_c\}\) and let
\(\widetilde H_{\Phi,\mathcal A_c}\) contain the active columns.  The covariance
of the active estimates is approximated by
\begin{equation}
\operatorname{Cov}(\widehat{\mathbf c}_{\mathcal A_c})
=
\widehat\sigma^2
(\widetilde H_{\Phi,\mathcal A_c}^\top
\widetilde H_{\Phi,\mathcal A_c}+\lambda I)^\dagger.
\label{eq:active_covariance}
\end{equation}
Here \(\widehat\sigma^2\) is a residual variance estimate that discounts the
effective degrees of freedom consumed by the fit; the use of
\(N-r_{\mathrm{eff}}(\tau_\sigma)\) is a heuristic degrees-of-freedom count that
does not exactly account for the background rank \(r\) or the NNLS active-set
constraints.  Equation~\ref{eq:active_covariance} is the standard active-set /
ridge least-squares covariance approximation \citep{seber2003linear}: its
diagonal entries give per-coefficient uncertainty and its off-diagonal entries
expose coupled (jointly uncertain) estimates.

For each transport ensemble member from
Section~\ref{subsec:wind_field_estimation}, IASA rebuilds
\(\widetilde H_\Phi^{(b)}\), refits \(\widehat{\mathbf c}^{(b)}\), and reports
empirical quantiles across members as transport uncertainty.  Every ensemble is
tagged by provenance.  Wind and dispersion members have
\texttt{ensemble\_kind=transport} and may form transport-uncertainty intervals.
Alternative inventory maps, source locations, labels, or scales have
\texttt{ensemble\_kind=inventory} and produce scenario-wise robustness tables.
The two kinds are never pooled into one interval.  Active-set intervals describe
the fitted observation model and likewise do not convert inventory scenarios
into probabilistic uncertainty.

\subsection{Residual Model-Adequacy Check}
\label{subsec:model_adequacy}

Residual magnitude alone is not calibrated evidence of misspecification.  Let
\(Z\) be an orthonormal basis for the background-orthogonal observed-row space,
and define the nondegenerate residual coordinates and externally supplied noise
covariance
\begin{equation}
\bar r=Z^\top(Y-H_\Phi^{\mathrm{lag}}\widehat{\mathbf c}),
\qquad
\bar\Sigma_e=Z^\top\Sigma_e Z.
\label{eq:adequacy_coordinates}
\end{equation}
where \(\Sigma_e\) is an externally declared observation-noise covariance and
\(Z\) removes the background directions so that \(\bar r\) carries no degrees of
freedom absorbed by \(Q\).  For independent homoscedastic noise
\(\Sigma_e=\sigma_e^2 I\), orthonormality of \(Z\) gives
\(\bar\Sigma_e=\sigma_e^2 I\) and \(T_{\mathrm{res}}=\|\bar r\|_2^2/\sigma_e^2\).
When \(\bar\Sigma_e\) is calibrated independently of this fitted residual and is
positive definite, the adequacy statistic is
\begin{equation}
T_{\mathrm{res}}=\bar r^\top\bar\Sigma_e^{-1}\bar r.
\label{eq:residual_adequacy_statistic}
\end{equation}
Its null distribution is not a plain \(\chi^2\), because \(\widehat{\mathbf c}\)
is fitted from the same data; IASA therefore calibrates it by parametric
bootstrap.  It simulates \(B_{\mathrm{res}}=1000\) synthetic observation sets from
the fitted source/background model plus the declared noise model, refits the
coefficients on each draw, and recomputes \(T_{\mathrm{res}}\) to form the null
distribution.  At \(\alpha=0.05\) it flags inadequacy when the observed statistic
exceeds the empirical \((1-\alpha)\)-quantile; the Monte Carlo \(p\)-value uses
the standard add-one correction.

If no externally calibrated noise model is supplied, IASA reports raw and
projected residual norms and sensor-wise, time-wise, and autocorrelation
summaries with \texttt{calibration\_status=uncalibrated}; it emits no pass/fail
adequacy claim.  Rejection establishes a mismatch with the declared model but
does not identify a missing source.  Non-rejection does not establish inventory
completeness: in particular, an omitted source whose sensor-time signature lies
in \(\operatorname{span}[H_\Phi^{\mathrm{lag}},Q]\) can be absorbed by fitted
terms and remain residual-invisible.

\subsection{Wind-Distribution Diagnostics}
\label{subsec:wind_distribution_diagnostics}

Diagnostics for one realized wind window answer a retrospective question.
Prospective network adequacy is assessed empirically over predeclared historical
or simulated wind windows.  Across those response matrices, IASA reports the
5th, 50th, and 95th percentiles of \(\sigma_J\), the probability of full
numerical and effective rank, the probability that each coefficient is weak,
the probability that each source pair is ambiguous, and the frequency of each
conservative report component.  These are Monte Carlo or empirical
distributional summaries under the sampled wind population, not a new
identifiability theorem.

\subsection{Acceptance Criteria for Refinement}
\label{subsec:refinement_acceptance}

If the future constrained refinement in
Section~\ref{subsec:constrained_refinement} is used, it is accepted only if it
improves fit without degrading separability.  Let
\(\widetilde H_{\Phi,0}\) be the projected response before refinement and
\(\widetilde H_{\Phi,\mathrm{ref}}\) after refinement.  We require
\begin{equation}
\sigma_J(\widetilde H_{\Phi,\mathrm{ref}})
\ge
(1-\eta_{\mathrm{id}})\sigma_J(\widetilde H_{\Phi,0})
\label{eq:refinement_smin_check}
\end{equation}
and
\begin{equation}
\max_{\substack{i\ne j\\i,j\notin\mathcal W}}\rho_{ij}^{\mathrm{ref}}
\le \tau_\rho^{\mathrm{ref}}.
\label{eq:refinement_coherence_check}
\end{equation}
These checks prevent end-to-end tuning from improving sensor fit by making
source--basis fingerprints less distinguishable.  If either response has fewer
than \(J\) numerically nonzero singular values, its \(\sigma_J\) is taken to be
zero (Section~\ref{subsec:identifiability_diagnostics}), so refinement cannot be
accepted by exploiting a rank-deficient response.

\subsection{Reporting Protocol}
\label{subsec:reporting_protocol}

The source ambiguity graph, its deterministic connected components, and the
transitive over-merging of an \(A\)--\(B\)--\(C\) chain are defined in
Section~\ref{subsec:source_merging}.  Operationally, those components are the
conservative recommended report groups, each edge retains its triggering
coefficient pair, and rank deficiency without an eligible pair yields a global
unresolved warning rather than an invented edge.

For each reported source group \(G_a\), IASA reports
\begin{align}
&\widehat{\boldsymbol\theta}_{G_a,1:T},
\quad \mathrm{CI}(\widehat\theta_{G_a}),
\quad \{v_{kb}:k\in G_a\},\nonumber\\
&\max_{\substack{k\in G_a,\,k'\notin G_a\\
                  (k,b),(k',b')\notin\mathcal W}}
\rho_{(k,b),(k',b')}.
\label{eq:reported_quantities}
\end{align}
A source contribution is marked reliable only when its constituent reported
coefficients are above the visibility threshold, its uncertainty interval is
sufficiently narrow, and its eligible coherence with other reported groups is
below the ambiguity threshold.  Reports retain the basis names of every weak
coefficient and the coefficient pair attaining each cross-source coherence
maximum.  Coherence and ray distance involving weak fingerprints are displayed
as ``N/A'' and excluded from extrema; the weak flags remain visible.  Otherwise
IASA labels the source as weakly visible, ambiguous with another source group,
or merged into a coarser identifiable group.

IASA additionally reports the per-sensor contribution decomposition and source
footprints of Section~\ref{subsec:per_sensor_footprints}, which show, for each
monitor, the fitted source-group contributions and their spatial cells of
origin.  These are spatial overlays of the same global fit: per-sensor source
shares are aggregated to the identifiable report groups \(\mathcal G\) wherever
the pooled diagnostics require grouping, so the spatial view never asserts a
separation the network does not support.
